\ChapterPage{Páginas 1--2 \textbullet{} La Gaceta}{Un despacho desde Greyhaven}
\begin{multicols}{2}
\section*{La campana bajo el rompeolas}
\textsc{Greyhaven, 3 de octubre.} Ayer sonó una campana bajo el puerto a las
tres y media, aunque el viejo rompeolas lleva sin campana desde la gran
tormenta de 1849. Tres pescadores la oyeron desde el muelle oriental. Cada
uno describió la nota como lenta, inmensa y, de algún modo, húmeda.

El capitán del puerto culpó a una boya suelta. Sin embargo, no faltaba boya
alguna, y el sonido subió por las suelas de las botas antes de alcanzar el
oído. Las gaviotas abandonaron los tejados en una sola nube blanca. Hasta
los perros de la calle Peregrina se negaron a cruzar la plaza.

\textit{Lorem abyssus}: en el latín del viejo impresor, el abismo lleva sus
propias cuentas. Cada ola anota un nombre; cada marea que se retira deja una
línea sin escribir. La frase es absurda, según el profesor Saye, que nunca
ha pasado un invierno en Greyhaven.

\section*{Aviso a los lectores}
La Crónica ruega a los visitantes que no raspen el vidrio verde del muro
marino. Los fragmentos no son un tesoro ni, como afirmó un escolar
entusiasmado, las escamas de un santo. Son vidrio corriente, suave como un
párpado y frío incluso junto a la estufa.

En el mercado de pescado ha aparecido una marca de tiza en la puerta de la
aduana clausurada. Parece una espiral con un brazo roto. El sereno la ha
borrado dos veces. Por la mañana, el signo vuelve, más pálido y orientado
hacia el mar.

\columnbreak
\section*{El parte meteorológico}
Presión descendente; viento del nordeste; mala visibilidad más allá del
arrecife. El capitán Morrow informa de un resplandor bajo la línea de agua,
azul verdoso y constante como una farola. La tripulación bajó una plomada.
Volvió tibia, con la cuerda cortada limpiamente por debajo del tercer nudo.

El doctor Elian Voss, de visita para catalogar pozas litorales, ha pedido al
periódico que corrija su crónica. Dice que la luz probablemente procede de
organismos fosforescentes. Lo dice mientras se niega a mirar el puerto
después del ocaso.

El cajista de la Crónica ha pedido el traslado. Asegura que los tipos de la
palabra \textsc{debajo} aparecen al revés en la forma de impresión. El
capataz lo atribuye a un descuido. El cajista lleva veinte años aquí y nunca
ha escrito mal un titular.

\section*{Editorial}
Se recomienda a los lectores mantener las lámparas preparadas y las puertas
cerradas. Si vuelve a sonar la campana, no cuenten sus golpes. Alguien lo
intentó en 1849, y el informe oficial termina a mitad de una frase.

\textbf{Tabla de mareas de mañana:} pleamar a las 6:12; bajamar a las 11:47.
El almanaque no incluye una marea que llegue desde tierra adentro.
\end{multicols}
